\usepackage[T1]{fontenc}
\usepackage{amsmath}
\usepackage{amsthm}
\usepackage{mathtools}
\usepackage{newpxtext}
\usepackage[vvarbb]{newpxmath}
\usepackage[varqu,varl,scaled=0.95]{zi4}
\usepackage{microtype}
\usepackage[a4paper,includehead,top=1cm,bottom=1.2cm,left=1.2cm,right=1.2cm,headheight=11pt,headsep=8pt]{geometry}
\usepackage{parskip}
\usepackage{xcolor}
\usepackage{xparse}
\usepackage{enumitem}
\usepackage{titlesec}
\usepackage{titletoc}
\usepackage{fancyhdr}
\usepackage{multicol}
\usepackage{listings}
\usepackage[most]{tcolorbox}
\usepackage[hypertexnames=false]{hyperref}
\usepackage{thmtools}

% Code copied out of the PDF comes back as plain ASCII (INC allows pasting
% from the notebook).
\input{glyphtounicode}
\pdfgentounicode=1

\newcommand{\teamname}{trieHard}
\newcommand{\university}{Institut Teknologi Bandung}

\definecolor{ink}{HTML}{1E293B}
\definecolor{muted}{HTML}{64748B}
\definecolor{thmcolor}{HTML}{1D4ED8}
\definecolor{defcolor}{HTML}{0F766E}
\definecolor{excolor}{HTML}{B45309}
\definecolor{warncolor}{HTML}{B91C1C}

\AtBeginDocument{\color{ink}}
\frenchspacing

\hypersetup{
  colorlinks=true,
  linkcolor=thmcolor,
  citecolor=thmcolor,
  urlcolor=thmcolor,
  pdftitle={\teamname\ ICPC Notebook},
  pdfauthor={\teamname, \university}
}

\setlength{\parskip}{3pt plus 1pt}
\setlist{itemsep=1pt,topsep=2pt,parsep=0pt,leftmargin=1.2em}
\setlist[enumerate,1]{label=\textup{(\arabic*)}}
\setcounter{secnumdepth}{2}
\setcounter{tocdepth}{2}

\makeatletter
\g@addto@macro\normalsize{%
  \setlength{\abovedisplayskip}{4pt plus 2pt minus 1pt}%
  \setlength{\belowdisplayskip}{4pt plus 2pt minus 1pt}%
  \setlength{\abovedisplayshortskip}{2pt plus 1pt}%
  \setlength{\belowdisplayshortskip}{2pt plus 1pt}%
}
\makeatother

\setlength{\columnsep}{14pt}
\setlength{\columnseprule}{0.4pt}
\renewcommand{\columnseprulecolor}{\color{muted!35}}
\raggedcolumns

\titleformat{\section}
  {\normalfont\large\bfseries}
  {\textcolor{thmcolor}{\thesection}}{0.6em}{}
  [\nobreak\vspace{1pt}{\color{thmcolor}\hrule height 0.8pt}\nobreak]
\titleformat{name=\section,numberless}
  {\normalfont\large\bfseries}
  {}{0pt}{}
  [\nobreak\vspace{1pt}{\color{thmcolor}\hrule height 0.8pt}\nobreak]
\titleformat{\subsection}
  {\normalfont\normalsize\bfseries}
  {\textcolor{thmcolor}{\thesubsection}}{0.6em}{}
\titleformat{\subsubsection}[runin]
  {\normalfont\normalsize\bfseries}
  {}{0pt}{}[.]
\titlespacing*{\section}{0pt}{10pt}{5pt}
\titlespacing*{\subsection}{0pt}{7pt}{3pt}
\titlespacing*{\subsubsection}{0pt}{4pt}{0.5em}

\titlecontents{section}[1.6em]
  {\addvspace{3pt}\bfseries}
  {\contentslabel[\textcolor{thmcolor}{\thecontentslabel}]{1.6em}}
  {\hspace*{-1.6em}}
  {\hfill\contentspage}
\titlecontents{subsection}[3.8em]
  {}
  {\contentslabel[\textcolor{muted}{\thecontentslabel}]{2.2em}}
  {\hspace*{-2.2em}}
  {\titlerule*[0.6em]{\textcolor{muted}{.}}\contentspage}

\pagestyle{fancy}
\fancyhf{}
\renewcommand{\headrulewidth}{0pt}
\renewcommand{\sectionmark}[1]{\markboth{#1}{}}
\renewcommand{\subsectionmark}[1]{}
\fancyhead[L]{\small\color{muted}\nouppercase{%
  \IfMarksEqualTF{2e-left}{first}{last}
    {\FirstMark{2e-left}}
    {\FirstMark{2e-left}\,--\,\LastMark{2e-left}}}}
\fancyhead[R]{\small\color{muted}\thepage}
% Page one. The rules ask for team and university at its top right.
\fancypagestyle{first}{%
  \fancyhf{}%
  \fancyhead[L]{\small\color{muted}ICPC Notebook}%
  \fancyhead[R]{\small\textbf{\teamname}\,\textperiodcentered\,\university
    \quad\textcolor{muted}{\thepage}}}

\newtheoremstyle{primerthm}{0pt}{0pt}{}{}{\bfseries\color{thmcolor}}{.}{0.5em}
  {\thmname{#1}\thmnote{\ \textnormal{(#3)}}}
\newtheoremstyle{primerdef}{0pt}{0pt}{}{}{\bfseries\color{defcolor}}{.}{0.5em}
  {\thmname{#1}\thmnote{\ \textnormal{(#3)}}}
\newtheoremstyle{primerex}{0pt}{0pt}{}{}{\bfseries\color{excolor}}{.}{0.5em}
  {\thmname{#1}\thmnote{\ \textnormal{(#3)}}}
\newtheoremstyle{primerrem}{}{}{}{}{\bfseries\itshape\color{muted}}{.}{0.5em}
  {\thmname{#1}\thmnote{\ \textnormal{(#3)}}}

\declaretheorem[name=Theorem,style=primerthm,numbered=no]{theorem}
\declaretheorem[name=Lemma,style=primerthm,numbered=no]{lemma}
\declaretheorem[name=Corollary,style=primerthm,numbered=no]{corollary}
\declaretheorem[name=Definition,style=primerdef,numbered=no]{definition}
\declaretheorem[name=Example,style=primerex,numbered=no]{example}
\declaretheorem[name=Remark,style=primerrem,numbered=no]{remark}

\tcbset{
  primerbox/.style={
    enhanced,
    breakable,
    lines before break=3,
    arc=0pt,
    outer arc=0pt,
    boxrule=0pt,
    leftrule=2pt,
    left=5pt,
    right=4pt,
    top=3pt,
    bottom=3pt,
    before skip=5pt,
    after skip=5pt
  }
}
\tcolorboxenvironment{theorem}{primerbox,colback=thmcolor!5,colframe=thmcolor}
\tcolorboxenvironment{lemma}{primerbox,colback=thmcolor!5,colframe=thmcolor}
\tcolorboxenvironment{corollary}{primerbox,colback=thmcolor!5,colframe=thmcolor}
\tcolorboxenvironment{definition}{primerbox,unbreakable,colback=defcolor!6,colframe=defcolor}
\tcolorboxenvironment{example}{primerbox,colback=excolor!4,colframe=excolor}
\newtcolorbox{notes}[1][Notes]{
  primerbox,
  colback=muted!8,
  colframe=muted,
  before upper={\textbf{\textcolor{muted}{#1}}\par
    \setlist[itemize]{leftmargin=1.1em,label=\textcolor{muted}{\textbullet}}}
}
\newtcolorbox{pitfall}[1][Pitfall]{
  primerbox,
  colback=warncolor!5,
  colframe=warncolor,
  before upper={\textbf{\textcolor{warncolor}{#1.}}\ }
}

\newcommand{\codesize}{\fontsize{7.5}{8.6}\selectfont}
\lstdefinestyle{cp}{
  language=C++,
  basicstyle=\ttfamily\codesize,
  keywordstyle=\color{thmcolor}\bfseries,
  commentstyle=\color{muted},
  stringstyle=\color{excolor},
  directivestyle=\color{defcolor},
  morekeywords={constexpr,nullptr,decltype,noexcept,static_assert,
    override,final,alignas,alignof,thread_local,
    ll,vi,pii,vector,pair,string,array},
  emphstyle=\color{defcolor},
  showstringspaces=false,
  columns=fullflexible,
  keepspaces=true,
  tabsize=2,
  upquote=true,
  breaklines=true,
  postbreak=\mbox{\textcolor{muted}{\(\hookrightarrow\)}\space},
  aboveskip=0pt,
  belowskip=0pt
}
\lstdefinestyle{cpinline}{style=cp,basicstyle=\ttfamily}
% Code boxes. title={gen.cpp} names the box, usage={gen 25 > in} adds the
% command to type, flushed right in the title bar (escape & as \&).
\tcbset{
  codebox/.style={
    primerbox,
    unbreakable,
    colback=ink!3,
    colframe=ink!35,
    colbacktitle=ink!10,
    coltitle=ink,
    fonttitle=\ttfamily\codesize\bfseries,
    toptitle=1.5pt,
    bottomtitle=1.5pt,
    listing only
  },
  usage/.style={after title={\hfill\mdseries\textcolor{muted}{#1}}},
  funcs/.style={listing options={style=cp,emph={#1}}}
}
\DeclareTCBListing{cpp}{!O{}}{
  codebox,
  listing options={style=cp},
  #1
}
\newtcbinputlisting{\cppfile}[2][]{
  codebox,
  listing file={#2},
  listing options={style=cp},
  #1
}
\lstdefinestyle{plain}{style=cp,language={}}
\newtcbinputlisting{\plainfile}[2][]{
  codebox,
  listing file={#2},
  listing options={style=plain},
  #1
}
\newcommand{\code}{\lstinline[style=cpinline]}

\NewDocumentCommand{\topic}{m !o}{%
  \subsection[#1]{#1\IfValueT{#2}{\hfill\textnormal{\textcolor{muted}{\(\bigO{#2}\)}}}}}

\newcommand{\N}{\mathbb{N}}
\newcommand{\Z}{\mathbb{Z}}
\newcommand{\Q}{\mathbb{Q}}
\newcommand{\R}{\mathbb{R}}
\newcommand{\C}{\mathbb{C}}
\newcommand{\bigO}[1]{\mathcal{O}(#1)}
\newcommand{\term}[1]{\textbf{\textcolor{defcolor}{#1}}}

\DeclareMathOperator{\lcm}{lcm}

\DeclarePairedDelimiter{\abs}{\lvert}{\rvert}
\DeclarePairedDelimiter{\floor}{\lfloor}{\rfloor}
\DeclarePairedDelimiter{\ceil}{\lceil}{\rceil}

% Writing kit
%   \topic{Dijkstra}[m \log n]      subsection with its complexity flushed right
%   \begin{cpp} ... \end{cpp}       C++ snippet, \begin{cpp}[title=...] for extra box options
%   \cppfile{code/dsu.cpp}          C++ snippet read from a file, paths relative to main.tex
%   \cppfile[breakable]{...}        lets a long snippet split across columns
%   \cppfile[title=gen.cpp,usage=gen 25 > in]{...}  title bar, usage on the right
%   \cppfile[funcs={update,query}]{...}  colors the names this snippet defines
%   \plainfile{code/x/build.txt}    same box for shell commands, no C++ colors
%   \code|vector<int>|              inline C++
%   theorem, lemma, corollary       blue boxes, \begin{theorem}[Fermat]
%   definition, example, remark     teal, amber and muted
%   \begin{notes}[Reading the ...]  grey box with a heading line, built for itemize
%   \begin{pitfall} ... \end{pitfall}  red box for bugs that cost you a submission
%   \bigO{n}, \floor{x}, \ceil{x}, \abs{x}, \lcm, \term{word}
%
% Limits measured for this layout: a code line holds 63 characters (safe
% even if the code font has to grow to 8pt), a full column holds about 87
% code lines, and the notebook may have at most 25 pages.
